%%%%%%%%%%%%%%%%%%%%%%%%%%%%%%%%%%%%%%%%%%%%%%%%%%%%%%%%%%%
% --------------------------------------------------------
% RMxAA Rho
% LaTeX Template
% Version 4.7 (10/08/2026)
%
% Authors: 
% Carlos Román-Zúñiga (croman@astro.unam.mx)
% Irene Cruz-González )icruzgonzalez@astro.unam.mx)
% Contributors:
% Guillermo Jimenez (memo.notess1@gmail.com)
% Eduardo Gracidas (eduardo.gracidas29@gmail.com)
% 
% License:
% XXXXXXXXXX
% --------------------------------------------------------
%%%%%%%%%%%%%%%%%%%%%%%%%%%%%%%%%%%%%%%%%%%%%%%%%%%%%%%%%%%
% --------------------------------------------------------
%                     ACKNOWLEDGMENTS
% This template has been developed based on the original rho 
% LaTeX class work, created by Luis Guillermo Jimenez Lopez 
% and Eduardo Gracidas Reyes. Available in the Overleaf's 
% template gallery.
% --------------------------------------------------------
%%%%%%%%%%%%%%%%%%%%%%%%%%%%%%%%%%%%%%%%%%%%%%%%%%%%%%%%%%%
% The following lines must be included in all contributions
%
\documentclass[10pt,letter,twoside]{rmat-rho}
\RMxAAtemplatetype{\RMxAT} % Select your article type
% {\RMxAA} article template
% {\RMxAC} conference template
\newcommand\rmaatex{RMxAT~\LaTeX}
\newcommand{\CS}[1]{\texttt{\textbackslash #1}}

%----------------------------------------------------------
% JOURNAL/HEADER INFORMATION
%----------------------------------------------------------

\vol{01}
% Counter page
\setcounter{page}{1000}
\pages{100-106}
\thisyear{2026}

%\doi{\href{https://www.astroscu.unam.mx/rmaa/RMxAA..XX-X}{\textcolor{blue}{https://www.astroscu.unam.mx/rmaa/RMxAA..XX-X}}}
\doi{\href{https://doi.org/10.22201/ia.XXXXXXXXXp.20XX.XX.XX.XX}{https://doi.org/10.22201/ia.XXXXXXXXp.20XX.XX.XX.XX}}
%\doi{ccccccccc}
\newenvironment{Example}
{\begin{list}{}{\setlength{\leftmargin}{10pt}\setlength{\rightmargin}{10pt}}%
  \item[]\itshape}
  {\end{list}}

%\usepackage{background}
\usepackage{lipsum}
\definecolor{textcolor}{HTML}{F9C051}

% From here on the RMxAC extended article, shortpaper or only abstract information should be included by the author(s)
%
%----------------------------------------------------------
% TITLE
%----------------------------------------------------------
%\journalname{\Large Revista Mexicana de Astronomía y Astrofísica 62, 1-10 (2026) \hfill {\includegraphics[width=1.2in, height=0.5in]{LogoRMxACamarillo.png}} \\}

\title{Thesis summary  for \rmaatex ~ template  v4.7}

%----------------------------------------------------------
% AUTHORS AND AFFILIATIONS
%---------------------------------------------------------

\author[1,$\dagger$]{A. Bright-Student,}
%\author[1,$\dagger$]{S. Upper-Visedby \orcidlink{0000-0001-2345-6789}}
%\author[2]{Segundo A. Cesor \orcidlink{0000-0009-8765-4321}}
%\author[1,2,$\dagger$]{A. Student}
%\author[2,4]{M.~Y.~Posdoc}

%----------------------------------------------------------

\affil[1]{Universidad Nacional Autónoma de México, Instituto de Astronomía, AP 70-264, CDMX 04510, México}
%\affil[2]{Universidad Nacional Autónoma de México, Instituto de Radioastronomía y Astrofísica.\\
%Antigua Carretera a Pátzcuaro 8701, Ex-Hda. San José de la Huerta, 58089, Morelia, Michoacán, México}

%\affil[4]{One more affiliation you may need}
%\affil[$\dagger$]{Posgrado en Astrofísica, UNAM}
%----------------------------------------------------------

\keywords{keyword 1, keyword 2, keyword 3, keyword 4, keyword 5}

%\pclave{palabra clave 1, palabra clave 2, palabra clave 3}
\received{\today}
%\revised{April 16, 2024}
\accepted{\today}
%\published{May 21, 2024}
%\doi{\url{https://www.astroscu.unam.mx/rmaa/RMxAA..XX-X}}
%----------------------------------------------------------
% DATES
%----------------------------------------------------------

%\dates{This manuscript was compiled on June 25th, 2024}

%----------------------------------------------------------
% FOOTER INFORMATION
%----------------------------------------------------------
\leadauthor{A. Bright-Student et al.}

%----------------------------------------------------------
% ARTICLE INFORMATION
%----------------------------------------------------------
% Correspondig author email
%---
\corres{studentemail@astro.unam.mx}
%\corres{Irene Cruz-González}
%\email{icruzgonzalez@astro.unam.mx}
%---

\received{\today}
%\revised{April 16, 2024}
\accepted{\today}
%\published{May 21, 2024}

%\license{Rho LaTeX Class \ccLogo\ This document is licensed under Creative Commons CC BY 4.0.}

% Include d Abstract and Resumen
%\setbool{thesislevel}{true} 

\setbool{rho-abstract}{false} % Set false to hide the abstract
\setbool{rho-resumen}{true} % Set false to hide the abstract
%\thesislevel{PhD}
%----------------------------------------------------------
% ABSTRACT
%----------------------------------------------------------

%\begin{thesislevel}
% prueba
%\end{thesislevel}

\begin{resumen}
Bienvenidos a los nuevos macros de la \rmaatex\  (v4.7 10/08/2026) para preparar contribuciones para ser publicados en la {\bf Revista Mexicana de Astronomía y Astrofísica Serie de Carrera Temprana (RMxAT)}. Utiliza la {\it LaTeX class}  {\bf rmat-rho.cls} desarrollada para la RMxAT. Hay dos tipos de macros de la Serie de Carrera Temprana:  RMxAT\_thesis y RMxAT\_memorias.% Archivos con ejemplos “RMxAT\_extended", "RMxAT\_shortpaper” or “RMxAT\_abstract ”, dependiendo del caso. 

El resumen en extenso de la tesis deberá ser presentado en inglés para que tenga una difusión global. Se supone que autor está familiarizado con elementos esenciales de LaTeX. 
%\\Este archivo {\it rmac\_extended\_example.tex} es el ejemplo de una contribución  {\bf en extenso}, destinado a las contribuciones orales invitadas y/o en extenso, que serán enviadas a los Editores Especiales del volumen correspondiente para ser publicadas en la RMxAC.

Para las tesis doctorales la extensión \textbf{máxima} del resumen es de \textbf{5} cuartillas. Para Maestría, será de \textbf{4} a lo máximo. Para Licenciatura el límite es de \textbf{3} cuartillas. En todos los casos ya incluyendo las figuras. Se recomienda de igual forma consultar los lineamientos de la revista para el envio de los documentos.

La RMxAT publicará también memorias de reuniones de astrónomos jóvenes utilizando el macro RMxAT\_memorias. Un ejemplo de este tipo de eventos es la Reunión de Estudiantes de Astronomía (REA) que se realiza anualmente en México. Sin embargo, eventos similares en otras partes de Iberoamérica pueden ser consideradas, previa consulta con la Jefatura Editorial de la RMxAT.
\end{resumen}




%----------------------------------------------------------

%\thesislevel{This is a PhD thesis realized for the Astrophysics graduate program at Universidad Nacional Autónoma de México}

\thesislevel{PhD}
\studiesprogram{Posgrado en Astrofísica}
\thesisinstitution{Universidad Nacional Autónoma de México}
\supervisors{S. Upper-Vissor (UNAM, Mexico) and Segundo A. Cesor (University of Moldavia)}

\begin{document}

\maketitle
\pagestyle{fancy}
\thispagestyle{firststyle}



%----------------------------------------------------------

\section{Introduction}

\RMxAAstart{T}here are two types of macros for contributions for publication in the Serie de Carrera Temprana (RMxAT): memorias and thesis. %Examples are in the files “rm\-extended\_memorias”, “rm\_shortpaper\_example” or “rm\_abstract\_example”, depending  
The style of all the templates are based on the \textit{rmat rho class} style. This requires minimal or no typesetting adjustments to provide a version of your manuscript that is close to the final printed version. This style also has ample margins to allow for a comfortable number of words per line and to leave room for marginal notes. %There is a complimentary author guide with details {\it RMxAA-authorguide}, which requires some familiarity with the use of \LaTeX{} and {\it Overleaf online editor} \url{https://www.overleaf.com}, where the RMxAC macros in \LaTeX{} can be obtained and used. 
For the author who requires a general introduction to \LaTeX{}, we recommend starting at {\it The LaTeX Project} website \url{https://www.latex-project.org/about/}, or using the Overleaf LaTeX guide\url{https://www.overleaf.com/learn/latex/Free_online_introduction_to_LaTeX_(part_1)}. 

\section{Preamble}

The first line to appear in your document should be 

\bigskip
\CS{documentclass}\verb+{rmat-rho}+

\bigskip
or

\bigskip
\CS{documentclass}[\textbf{optionlist}]\verb+{rmat-rho}+

\bigskip
\noindent which sets up the document to use the RMxAT class, using the default \texttt{manuscript} option, which is designed for use by authors who submit extended articles to the RMxAT journal. 

The following commands can be used after the \CS{documentclass} command and before the command \CS{begin}\verb|{document}|: %{\mathrm{document}\}$:

\begin{itemize}
    \item \CS{title} -- defines article title
    \item \CS{author} command defines the authors of the article.
    \item \CS{affil} -- defines the author´s affiliations
    \item \CS{leadauthor} -- provides the last name and initials of the leading author of the article, and will be visible at the top of every odd-number page.
    \item \CS{smalltitle} -- defines the footer that appears at the bottom of the page of every article in RMxAC, and it is used to register the details of the meeting: conference title, city and country, date, and special volume editors.
    \item \CS{corres} and \CS{email} define the name and email address of the corresponding author. Note that you can use \CS{corres} or \CS{email} to list only the email address of the corresponding author.
\end{itemize}

\section{Keywords}

A minimum of three and a maximum of five keywords must be provided using the command \verb|\keywords{keyword1, keyword2, etc.}|. 


\section{Abstract}

The {\it abstract} and {\it resumen} of the article are placed between the commands: \CS{begin}\verb|{abstract}| ... \CS{end}\verb|{abstract}| and \CS{begin}\verb|{resumen}| ... \CS{end}\verb|{resumen}|. The recommended size for both is a maximum of 200 words, and they should not exceed 300 words.

If the authors are unable to provide a Spanish version of the abstract, they can use the same text as the English abstract, and our editors will take care of the Spanish version.
 



\section{Main Body}

The main body of the document should be enclosed within the following pair of commands

\bigskip
\CS{begin}\verb+{document}+

...

ARTICLE TEXT

...

\CS{end}\verb+{document}+
\bigskip

The first command after the \CS{begin}\verb+{document}+ should be \CS{maketitle}. This formats the title, author names, abstract, and keywords.

Within the main body of the document, all standard \LaTeX{} commands can be used. The commands provided by the many optional packages distributed with \LaTeX{} may also be used as long as the package is loaded using the \CS{usepackage} command in the preamble. However, authors are requested to avoid using commands that change document fonts, page layout, or other ``stylistic'' parameters. We should also note that not all optional packages have been tested for compatibility with the rmaa\_rho class.

\subsection{Sectioning commands and cross references}

Authors are encouraged to use the standard \LaTeX{} sectioning commands to subdivide their articles:

\bigskip

\CS{section}

\CS{subsection}

\CS{subsubsection}

\CS{paragraph}

\bigskip

Please use standard \LaTeX{} sectioning commands to subdivide your document into appropriate sections. You should use mixed cases for the section titles, although in the current style, this only really matters at the level of \CS{subsection} and below.

These will be automatically typed in the RMxAA and RMxAC styles, respectively. Cross-referencing is made easier by the use of the \CS{label} 
 \verb+{LABEL}+ command immediately after each sectioning command, where the LABEL text is any mnemonic string. Elsewhere in the document, the section can then be referred to as \S \CS{ref} \verb+{LABEL}+. The \CS{label} command can also be used with equations, figures, and tables (see below). 

It is preferable to use the \CS{label}/\CS{ref} mechanism for cross references in order to minimize the chance of errors and
allow automatic hyperlinks in the PDF output. The style that should be used for cross references is, for example, Figure 3, Table 1, equation (12), and \S 5.1, where the symbol of the section “§” is produced by the \LaTeX{} command “\S”.

\subsection{Math Symbols and Equations}
\label{sec:math}

Symbols for physical quantities should usually be in italics: velocity,
$v$, density, $N$, etc. However, multi-letter symbols generally look
better in roman: FWHM, EM, etc. 

The subscripts should be in roman (coded using \CS{mathrm}) unless they are variables:
$N_\mathrm{e}$, $T_\mathrm{eff}$, but $\sum_i a_i$. 

Physical units should be in roman with thin spaces: $10\,\mathrm{K}$, $1.2\times
10^{-12} \,\mathrm{erg\,cm^{-2}\,s^{-1}}$, etc. 

Things generally come
best if you place an entire expression within a single pair of\$'s and then make judicious use of \CS{mathrm}. For example,
\begin{Example}
  $ \mathrm{FWHM} = \int N_\mathrm{e} N_\mathrm{i} \, dz $ 
\end{Example}

Recall that the ``minus sign" only exists inside math mode: minus two is $-2$, not~-2, nor \hbox{even --2!} In addition, remember that spacing inside the math mode is designed for equations, not words; therefore, you should not use\$ to obtain italic text. Compare $eff$ and \textit{eff}. 

The \CS{frac} command (and its \TeX{} relative \CS{over}) is best used only in displayed equations. Something like 
\begin{equation}
  \label{eq:one}
  x = \frac { a + b } { c } 
\end{equation}
looks fine, whereas $x = \frac { a + b } { c }$ looks somewhat cramped. Better rewritten as $x = (a + b) / c $. 

\paragraph{How to define a macro that can be used inside or outside  math mode.} Use the \CS{ensuremath} command. For example: 
\begin{verbatim}
\newcommand{\fluxunits}{%
   \ensuremath{\mathrm{%
      erg\,s^{-1}\,cm^{-2}}}}
\end{verbatim}
Then you can write either \verb+15.1\,\fluxunits+ or
\verb+$2.3\times+ \verb+10^{-11} \,+ \verb+\fluxunits$+

\subsubsection{Equations}

Equation~(\ref{ec:equation}), shows the Schrödinger equation as an example:     \begin{equation} \label{ec:equation}
         \frac{\hbar^2}{2m}\nabla^2\Psi + V(\mathbf{r})\Psi = -i\hbar \frac{\partial\Psi}{\partial t}
    \end{equation} 
Punctuation can be added as in Eq.~(\ref{ec:equation3})
    \begin{equation} \label{ec:equation3}
         \frac{\hbar^2}{2m} \nabla^2\Psi + V(\mathbf{r})\Psi = -i\hbar \frac{\partial\Psi}{\partial t}.
    \end{equation} 
    
\subsection{Special symbols}

The commands for the commonly used symbols in astronomy are presented in Table~\ref{tab:anysymbols}.

\begin{table}[ht]
\caption{Astronomy special symbols commands.}%Additional commands for special symbols commonly used in astronomy. These can be used anywhere.}
 \label{tab:anysymbols}
 \centering
 \begin{tabular*}{0.9\columnwidth}{@{}l@{\hspace*{12pt}}l@{\hspace*{12pt}}l@{}}
  \hline
  Command & Output & Meaning\\
  \hline
  \verb'\sun' & \sun & Sun, solar\\[2pt] % additional height spacing for enhanced symbol legibility
  \verb'\earth' & \earth & Earth, terrestrial\\[2pt]
  \verb'\micron' & \micron & microns\\[2pt]
  \verb'\degr' & \degr & degrees\\[2pt]
  \verb'\arcmin' & \arcmin & arcminutes\\[2pt]
  \verb'\arcsec' & \arcsec & arcseconds\\[2pt]
  \verb'\fdg' & \fdg & fraction of a degree\\[2pt]
  \verb'\farcm' & \farcm & fraction of an arcmin\\[2pt]
  \verb'\farcs' & \farcs & fraction of an arcsec\\[2pt]
  \verb'\fd' & \fd & fraction of a day\\[2pt]
  \verb'\fh' & \fh & fraction of an hour\\[2pt]
  \verb'\fm' & \fm & fraction of a minute\\[2pt]
  \verb'\fs' & \fs & fraction of a second\\[2pt]
  \verb'\fp' & \fp & fraction of a period\\[2pt]
 % \verb'\diameter' & \diameter & diameter\\[2pt]
  \verb'\sq' & \sq & square, Q.E.D.\\[2pt]
  \hline
 \end{tabular*}
\end{table}



\subsection{Ionization states}

For ions, a \verb'\ion{}{}' command is used for the correct typesetting of ionization states.
For example, to typeset singly ionized calcium, use \verb'\ion{Ca}{ii}', which produces \ion{Ca}{ii}, while a doubly ionized forbidden oxygen line produces [\ion{O}{iii}].

\section{Codes}

    This class\footnote{Hello there I am a footnote:)} includes the \textit{listings} package, which offers customized features for adding codes, especially for C, C++, \LaTeX, and Matlab. The format can be customized in the \textit{rho class} file.

%    \nolinenumbers
%    \lstinputlisting[caption=Example of matlab code., label={lst:listing-Mat}, language=Matlab]{example.m}
 %   \linenumbers
    
    If line numbering is enabled, we recommend placing the command \verb|\nolinenumbers| at the beginning and \verb|\linenumbers| at the end of the code, as follows: 
    
    This temporarily removes the line numbering, and the code looks better.
    


\section{Figures and Tables}

    \subsection{Sample figure}

        Figure~\ref{fig:figure} shows an example. The labels in each figure are recommended to use the usual fig:figurename.
        
            \begin{figure}[ht]
                \centering
                \includegraphics[width=0.6\columnwidth]{figures/example-fig.pdf}
                \caption{Example figure of size = 0.6 columnwith.}
                \label{fig:figure}
            \end{figure}

    \subsection{Sample double figure}

        Figure~\ref{fig:examplefloat} shows an example of a floating figure with two pictures that cover the width of the page. It can be placed at the top or the bottom of the page. The space between the figures can also be changed using the \verb|\hspace{Xpt}| command.

        \begin{figure*}[ht] % t for position at the top of the current page; b for position at the bottom of the current page
            \centering
                \begin{subfigure}[b]{0.38\linewidth} % Fig (a)
                    \includegraphics[width=0.95\linewidth]{figures/example2.pdf}
                    \caption{Example left figure.}
                    \label{fig:figa}
                \end{subfigure}
            \hspace{15pt}   % Space between the figures
                \begin{subfigure}[b]{0.38\linewidth} % Fig (b)
                    \centering
                    \includegraphics[width=0.95\linewidth]{figures/example2.pdf}
                    \caption{Example right figure.}
                    \label{fig:figb}
                \end{subfigure}
            \caption{Example figure that covers 0.9 of the width of the page obtained from PGFPlots \cite{PFGPlots}.}
            \label{fig:examplefloat}
        \end{figure*}

    \subsection{Sample table}

    Similar to figures, tables can be placed in one or two columns, depending on their length.

        Table~\ref{tab:table2}, shows an example table that covers the width of the page positioned at the bottom of a new page. Table~\ref{tab:table2}, shows an example of a column-width table positioned at the top of a page. The labels in each table are recommended using the usual tab:tablename.

        \begin{table*}[ht]
            \RaggedRight
            \caption{Table example that covers the width of the page.}
            \label{tab:table2}
                \begin{tabular}{lllp{9.7cm}}
                    \toprule
                    \textbf{Day} & \textbf{Min Temp} & \textbf{Max Temp} & \textbf{Summary} \\ 
                    \midrule
                    Monday & 11\textdegree C & 22\textdegree C & A clear day with lots of sunshine.  Strong breezes lower the temperatures. \\
                    Tuesday & 9\textdegree C & 19\textdegree C & Cloudy with rain across many northern regions. \\
                    Wednesday & 10\textdegree C & 21\textdegree C & Rain will still linger in the morning. 
                    Conditions will improve by early afternoon and continue 
                    throughout the evening.\\
                    \bottomrule
                \end{tabular}
                
            \tabletext{Note: Obtained from Latex tables \cite{projects-2023}.}
            
        \end{table*}

        
        \begin{table}[ht]
            \RaggedRight
            \caption{Table example that covers one column of the page.}
            \label{tab:table3}
                \begin{tabular}{lccl}
                    \toprule
                    \textbf{Day} & \textbf{$T_{min}$} & \textbf{$T_{max}$} & Reference \\ 
                    \midrule
                    Monday & 11\textdegree C & 22\textdegree C & {\small \citet{1987flme.book.....L}}\\
                    Tuesday & 9\textdegree C & 19\textdegree C & {\small \citet{1987flme.book.....L}} \\
                    Friday & 10\textdegree C & 21\textdegree C & {\small \citet{1987flme.book.....L}}\\
                    \bottomrule
                \end{tabular}
                
            \tabletext{Note: Similar a Table~2, de ancho de una columna, con referencia y nota}.
            
        \end{table}

        If the table exceeds the width of the page, the environment \verb'\begin{landscape}' ... \verb'\end{landscape}' can be used. Unfortunately, this method forces a page break before the table. Authors should not worry, as this will be corrected during the editorial work.

\section{Appendices}

If you have appendices to the article, you can use something like the following: \hfill \\

\CS{begin}\verb|{appendices}|\\
\CS{section}\verb|{First Appendix}|\\
\CS{label}\verb|{sec:ap-A}}|\\
\CS{}\verb|{Text of first appendix.}|\\
\CS{section}\verb|{Second Appendix}|\\
\CS{label}\verb|{sec:ap-B}|\\
\CS{}\verb|{Text of second appendix.}|\\
\CS{end}\verb|{appendices}|\\

The appendices go after the acknowledgments but before the
bibliography. Equations in the appendices will be labeled
A1, A2, B1, B2, etc.

\section{Concept highlight box}\label{box:box1}

 
The new \rmaatex macro allows the authors to use a colored box to highlight a concept or equation, as shown in this example. The label and reference point to the section that is included. Example: See the Concept box in Section~\ref{box:box1}.  
    
        \begin{rhoenv}[frametitle=Highlight Concept Box,backgroundcolor=\rhotheme!22]
            Hello! I am an example of a concept highlight box. \label{box:boxy}
        \end{rhoenv}
    


%\section*{Unnumbered section} \label{sec:unsec}

%    If an unnumbered section is declared, a square appears followed by the section name. This style is characteristic of this class and is only applicable to first-level sections.

%    Since this affects the title of the table of contents and references, you can make a modification in \textit{section style-rho class} to remove the square. See the Appendix for more information.

\section{Reference style}

     Our default formatting for references uses the journal-naming system from the AASTex style and Astrophysics Data System (ADS). RMxAA journals follow the ADS bibliographic codes for both refereed \url{https://adsabs.harvard.edu/abs_doc/refereed.html} and non-refereed \url{https://adsabs.harvard.edu/abs_doc/non_refereed.html} publications.
    
    If a paper has more than five authors, only the first three will be listed, followed by ``et al.". The DOI of the publication was also included. At the end of this document, you will find an example of default reference formatting. 

    The default format for references follows the ADS BibTeX style. The author should provide a bibliography file using the command \CS{bibliography}\verb|{bibfile}|, where all references are included. At the end of this document, you will find an example of the default reference formatting. Note that the DOI code is linked to references that have one; therefore, it is important to include it in the bib file.
    
    The usual commands \CS{citep} for reference in parentheses, as in \citep{Roman23},  and \CS{citet} for references with year in parentheses, as in \citet{Roman23}, are used in the main body of the article.   
    
The {\bf rmxat.bib file} is an example of a bib file that should also be submitted with the manuscript.

%-------------References bib file------------------------

\bibliography{RMxAT}

%----------------------------------------------------------

\end{document}
